\section{Lifted formulas and comparison of local data}
\label{subsec:rg-local-pullback}

The extraction, adjustment and scaling maps of
Sections~\ref{sec:extraction} and~\ref{subsec:rg-scaling}
have the following local formulas, as in
\cite[\S12.9]{QED3Manuscript}. For comparison of different torus volumes we
need formulas that depend on a lift and local data alone.
The required margin also distinguishes the pairs produced
by scaling from the full set of coarse pairs.

For a finite photon set $\widetilde U$ on the covering
lattice, let $\mathfrak G_k(\widetilde U)$ be the algebra
generated by the ordinary indices over $\widetilde U$
and pair indices whose two photon sites lie in
$\widetilde U$. Its ordinary and pair measures remain
$\eta_k^3/48$ and $(\eta_k^3/48)^2$. For a system of
separately antisymmetric kernels $K=(K_{nm})$ define
\[
 \operatorname{El}_k(K)
   =\sum_{n,m}\frac1{n!m!}\int\dd\xi\,\dd\zeta\,
                 K_{nm}(\xi;\zeta)\psi(\xi_1)\cdots\psi(\xi_n)
                                    \chi(\zeta_1)\cdots\chi(\zeta_m).
\]
The integrals range over these two finite index sets.
The kernel expansion is a bijection between the systems
of such kernels and this algebra: at increasing arguments
its coefficient is the kernel times the positive product
of index weights, as in \eqref{eq:rg-activity-kernels}.

Map each covering generator to the torus generator at
its simultaneous period class, with the antiperiodic
convention, and denote the homomorphism by
$\operatorname{cov}_\ell$. If
\begin{equation}\label{eq:rg-generator-pullback-margin}
 \operatorname{diam}_\infty(\widetilde U)+\eta_k<\ell,
\end{equation}
it is injective and preserves kernels when read on this
lift. In fact two congruent photon sites in
$\widetilde U$ must coincide. Choose a fundamental
domain containing $\widetilde U$ and order its
independent ordinary and pair indices so that they
extend the covering orders. The homomorphism maps
distinct covering basis monomials to distinct torus
basis monomials with the same coefficients and weights.
This proves injectivity, preservation of bidegrees,
and equality of the kernels at the lifted arguments.
All torus kernels at indices outside these period
classes vanish.

Now let $\widetilde X$ be a face-connected lift and put
$\widetilde U=\widetilde X^{+\eta_k}$, its mesh
neighborhood. The margin
\begin{equation}\label{eq:rg-activity-pullback-margin}
 \operatorname{diam}_\infty(\widetilde X)+2\eta_k+M<\ell
\end{equation}
gives a bijection of ordinary indices over
$\widetilde X$ onto those over $X$, and of admissible
pairs for $\widetilde X$ onto those for $X$.
For a pair of the first kind, its endpoints and every
point of its shortest paths lift inside $\widetilde X$:
a lifted path point within distance $M$ of its first
endpoint differs from its representative in
$\widetilde X$ by a period of norm at most
$M+\operatorname{diam}_\infty(\widetilde X)<\ell$.
That period is zero. For a boundary pair
$x'=x^{(3)}$, lift the endpoint in $X$; the other is
its prescribed nearest neighbor in $\widetilde U$.
Two lifted pairs in the same class coincide because
their first photon sites coincide.
These arguments prove surjectivity, injectivity and
preservation of pair kind.
The margin also implies
\eqref{eq:rg-generator-pullback-margin} for $\widetilde U$.
By the nonwrapping lifting result, $X$ is small exactly
when $d_M(\widetilde X)\le L$: for either a small torus
polymer or a small covering polymer the unique
nonwrapping lift has the same net length.

Take data $(K^0,K,\widetilde B)$ consisting of two
kernel systems over $\widetilde U$ and a connection
on bonds with both ends in $\widetilde X$.
For $d_M(\widetilde X)\le L$, define
\begin{equation}\label{eq:rg-lifted-extraction}
 \begin{split}
 M^0&=\operatorname{Vol}(\widetilde X)^{-1}
      \int_{\widetilde X}\dd x'\int_{\widetilde X}\dd x''\,
                                                    K^0_{20}(x',x''),\\
 \Phi^R&=\operatorname{El}_k(K)-K^0_{00}
               -\int_{\widetilde X}\dd x\,\bar\psi(x)^TM^0\psi(x),\\
 \Phi^E&=\Phi^R-(\Phi^R)^{[2,0]}_{\widetilde X}
               +\Theta_{\widetilde X,\widetilde B}
                                      (\Phi^R)^{[2,0]}_{\widetilde X}.
 \end{split}
\end{equation}
Here $K^0_{20}$ is the balanced matrix component in the
earlier convention, and the subscript $\widetilde X$
retains only quadratic terms whose two sites lie over
$\widetilde X$. Thus the transport homomorphism is
applied within its domain. For $d_M(\widetilde X)>L$,
set $\Phi^R=\Phi^E=\operatorname{El}_k(K)$. In both
cases define
\[
 \Phi^S=\widetilde S_L\Phi^E,
\]
where $\widetilde S_L$ is the covering version of
\eqref{eq:rg-leg-scaling}. These formulas use no side
$\ell$. Translation by $v\in M\Z^3$ preserves the
cube geometry, measures and chain transports. Therefore
\begin{equation}\label{eq:rg-lifted-translation}
 \Phi^R(T_v\mathrm{data})=T_v\Phi^R(\mathrm{data}),\qquad
 \Phi^E(T_v\mathrm{data})=T_v\Phi^E(\mathrm{data}),\qquad
 \Phi^S(T_v\mathrm{data})=T_{v/L}\Phi^S(\mathrm{data}),
\end{equation}
with the translated polymer understood on the left.

For an activity $\mathscr E$, let $K^0$ and $K$ be its
kernels at $(0,0)$ and at $(\mathscr A_L,\mathscr C_L)$,
restricted to the lift, and let $\widetilde B$ be the
restriction of $(\mathscr A+\mathscr C)_L$ to the
internal bonds. Then
\begin{equation}\label{eq:rg-lifted-activity-identities}
 \begin{split}
 (R\mathscr E)(X;\mathscr A_L,\mathscr C_L)
                              &=\operatorname{cov}_\ell\Phi^R,\\
 (R\mathscr E)^{\mathrm{nat}}(X;\mathscr A_L,\mathscr C_L)
                              &=\operatorname{cov}_\ell\Phi^E,\\
 S_L[(R\mathscr E)^{\mathrm{nat}}(X;\mathscr A_L,\mathscr C_L)]
                              &=\operatorname{cov}_{\ell/L}\Phi^S.
 \end{split}
\end{equation}
To prove the first identity, locality and the index
bijection reproduce the original activity from $K$.
The vacuum subtraction is exactly $K^0_{00}$, and
the two label-averaged mass integrals have the same
weights and volume in the lift and on the torus.
The diagonal bilinear is period invariant, since its
two antiperiodic signs cancel. The subtraction is
therefore also the same.
For the second identity the ordinary quadratic terms
of $\Phi^R$ lie over $\widetilde X$. Cubes, their
shortest chains and their transports are the lifted
versions of the torus construction. On each ordinary
generator this gives
$\operatorname{cov}_\ell\Theta_{\widetilde X,\widetilde B}
=\Theta_{X,(\mathscr A+\mathscr C)_L}\operatorname{cov}_\ell$.
The homomorphism identity proves the adjusted formula.
Finally, on each generator
\[
 \operatorname{cov}_{\ell/L}\widetilde S_L
                                  =S_L\operatorname{cov}_\ell.
\]
Both sides have the same scaled index and coefficient,
with compatible antiperiodic signs. This proves the
third identity.

The margin \eqref{eq:rg-activity-pullback-margin}
suffices for the image of $\widetilde S_L$, whose
pairs have $d_X\le M/L$. To identify
$\operatorname{cov}_{\ell/L}$ with the quotient map
on the full coarse local algebra, including pairs
with $M/L<d_X\le M$, we use the stronger margin
\begin{equation}\label{eq:rg-full-coarse-pullback-margin}
 \operatorname{diam}_\infty(\widetilde U)+LM<\ell.
\end{equation}
It is equivalent to
$\operatorname{diam}_\infty(\widetilde U/L)+M<\ell/L$.
For a small polymer both margins hold, since
\[
 \operatorname{diam}_\infty(\widetilde X)+2\eta_k+LM
                    \le2LM+4M<\ell.
\]
The torus expressions in
\eqref{eq:rg-lifted-activity-identities} do not depend
on the chosen lift. Translation of that lift by $\ell w$
translates the data; their total even degree makes the
product of their argument signs one.
Equation \eqref{eq:rg-lifted-translation} and the same
evenness then cancel the generator signs under the
quotient maps at both levels.
Thus identical lifted data produce identical lifted
extraction, adjustment and scaled kernels on different
admissible tori, subject to the stated margins.
